\documentclass[11pt,a4paper]{article}

% --- Packages ---
\usepackage[margin=.75in]{geometry}
\usepackage[utf8]{inputenc}
\usepackage[T1]{fontenc}
\usepackage{lmodern} % Fix for font not found error
\usepackage{titlesec}
\usepackage{enumitem}
\usepackage{hyperref}
\usepackage{xcolor}
\usepackage{fancyhdr}
\usepackage{comment}
\usepackage{cleveref}
\usepackage[bottom]{footmisc}
\usepackage{graphicx}
\usepackage{longtable}
\usepackage{booktabs}
\usepackage{array}
\usepackage{calc}
\usepackage{etoolbox}

% Standard Pandoc setup
\providecommand{\tightlist}{%
  \setlength{\itemsep}{0pt}\setlength{\parskip}{0pt}}

\crefformat{section}{\S#2#1#3}
\crefformat{subsection}{\S#2#1#3}
\crefformat{subsubsection}{\S#2#1#3}
\crefformat{figure}{#2Figure~#1#3}
\crefformat{footnote}{#2\footnotemark[#1]#3}

\linespread{0.9} % Riduce lo spazio tra le linee

% --- Colors ---
\definecolor{primary}{RGB}{0, 0, 0}
\definecolor{links}{HTML}{0000EE}
\definecolor{blue}{HTML}{26428B}

\hypersetup{
    colorlinks=true,
    linkcolor=links,
    urlcolor=links,
    citecolor=links,
}

% --- Section Formatting ---
\titleformat{\section}
  {\normalfont\Large\bfseries\color{blue}}{\thesection}{1em}{}
\titleformat{\subsection}
{\normalfont\large\bfseries\color{blue}}{\thesubsection}{1em}{}
\titleformat{\subsubsection}
{\normalfont\bfseries\color{blue}}{\thesubsubsection}{1em}{}

% --- Header and Footer ---
\pagestyle{fancy}
\fancyhead[L]{\small Giuseppe Capasso}
\fancyhead[R]{\small vapt}
\fancyfoot[C]{\thepage}
\renewcommand{\headrulewidth}{0.4pt}

% --- Custom Title Command ---
\newcommand{\proposaltitle}[1]{
    \begin{center}
        \vspace*{0.1cm}
        {\LARGE \bfseries \color{blue} #1} \\[1.5ex]
        {\large \textbf{Giuseppe Capasso}} \\[1ex]
                \small{
            \vspace{0.1cm}
            \textbf{Materia:} vapt\\
        }
            \end{center}
    \vspace{0.3cm}
    \hrule height 0.5pt
    \vspace{0.5cm}
}

\begin{document}

\proposaltitle{Nome e Cognome: _____________________________________________________}

Nome e Cognome:
\_\_\_\_\_\_\_\_\_\_\_\_\_\_\_\_\_\_\_\_\_\_\_\_\_\_\_\_\_\_\_\_\_\_\_\_\_\_\_\_\_\_\_\_\_\_\_\_\_\_\_\_\_
Data: \_\_\_\_\_\_\_\_\_\_\_\_\_\_\_\_\_\_\_

Punteggio: \_\_\_\_\_/20

\begin{enumerate}
\def\labelenumi{\arabic{enumi}.}
\tightlist
\item
  Qual è l'obiettivo principale della Binary Exploitation?
\end{enumerate}

\begin{enumerate}
\def\labelenumi{\Alph{enumi})}
\tightlist
\item
  Alterare il codice sorgente originale del programma
\item
  Aggiornare le difese del sistema operativo
\item
  Manipolare il flusso di esecuzione tramite input malformati per
  ottenere l'Arbitrary Code Execution
\item
  Crittografare i file sensibili sul server bersaglio
\end{enumerate}

\begin{enumerate}
\def\labelenumi{\arabic{enumi}.}
\setcounter{enumi}{1}
\tightlist
\item
  Nella memoria virtuale di un processo, quale area è specificamente
  dedicata all'allocazione dinamica (ad esempio quando si utilizzano
  funzioni come malloc())?
\end{enumerate}

\begin{enumerate}
\def\labelenumi{\Alph{enumi})}
\tightlist
\item
  Stack
\item
  Heap
\item
  Text
\item
  BSS
\end{enumerate}

\begin{enumerate}
\def\labelenumi{\arabic{enumi}.}
\setcounter{enumi}{2}
\tightlist
\item
  Per quale motivo strutturale il linguaggio C è particolarmente esposto
  alla vulnerabilità di Buffer Overflow?
\end{enumerate}

\begin{enumerate}
\def\labelenumi{\Alph{enumi})}
\tightlist
\item
  Perché la memoria viene gestita da un Garbage Collector troppo lento
\item
  Perché non effettua alcun controllo automatico dei limiti (bounds
  checking) sulle allocazioni
\item
  Perché richiede permessi di root per ogni esecuzione
\item
  Perché non supporta l'utilizzo dei registri della CPU
\end{enumerate}

\begin{enumerate}
\def\labelenumi{\arabic{enumi}.}
\setcounter{enumi}{3}
\tightlist
\item
  Qual è la dimensione standard di un Indirizzo Fisico (MAC Address)?
\end{enumerate}

\begin{enumerate}
\def\labelenumi{\Alph{enumi})}
\tightlist
\item
  32 bit
\item
  48 bit
\item
  64 bit
\item
  128 bit
\end{enumerate}

\begin{enumerate}
\def\labelenumi{\arabic{enumi}.}
\setcounter{enumi}{4}
\tightlist
\item
  Come si comporta uno switch di rete quando riceve un frame con un MAC
  Address di destinazione non presente nella sua tabella (Unknown
  Unicast)?
\end{enumerate}

\begin{enumerate}
\def\labelenumi{\Alph{enumi})}
\tightlist
\item
  Scarta immediatamente il frame per sicurezza
\item
  Rimanda il frame al mittente con un messaggio di errore
\end{enumerate}

Nome e Cognome:
\_\_\_\_\_\_\_\_\_\_\_\_\_\_\_\_\_\_\_\_\_\_\_\_\_\_\_\_\_\_\_\_\_\_\_\_\_\_\_\_\_\_\_\_\_\_\_\_\_\_\_\_\_
Data: \_\_\_\_\_\_\_\_\_\_\_\_\_\_\_\_\_\_\_

Punteggio: \_\_\_\_\_/20

\begin{enumerate}
\def\labelenumi{\Alph{enumi})}
\setcounter{enumi}{2}
\tightlist
\item
  Effettua il Flooding, inoltrando il frame su tutte le porte attive
  tranne quella di origine
\item
  Mette il frame in coda aspettando che il destinatario si presenti
\end{enumerate}

\begin{enumerate}
\def\labelenumi{\arabic{enumi}.}
\setcounter{enumi}{5}
\tightlist
\item
  Quale debolezza architetturale fondamentale del protocollo ARP viene
  sfruttata nell'attacco di ARP Spoofing?
\end{enumerate}

\begin{enumerate}
\def\labelenumi{\Alph{enumi})}
\tightlist
\item
  La mancanza di qualsiasi meccanismo di autenticazione per i messaggi
  inviati
\item
  L'uso di chiavi crittografiche deboli a 56-bit
\item
  La lentezza intrinseca nella propagazione dei pacchetti broadcast
\item
  La frammentazione dei pacchetti a livello 3
\end{enumerate}

\begin{enumerate}
\def\labelenumi{\arabic{enumi}.}
\setcounter{enumi}{6}
\tightlist
\item
  Qual è uno dei principali impatti di un attacco di ARP Spoofing (o
  Poisoning) riuscito?
\end{enumerate}

\begin{enumerate}
\def\labelenumi{\Alph{enumi})}
\tightlist
\item
  La corruzione del firmware del router
\item
  L'esecuzione di uno shellcode remoto
\item
  L'intercettazione del traffico tramite una posizione di
  Man-in-the-Middle (MitM)
\item
  La cancellazione della MAC Address Table dello switch
\end{enumerate}

\begin{enumerate}
\def\labelenumi{\arabic{enumi}.}
\setcounter{enumi}{7}
\tightlist
\item
  Cosa può accadere se un attaccante riesce a sovrascrivere illegalmente
  il ``Return Address'' salvato nello Stack?
\end{enumerate}

\begin{enumerate}
\def\labelenumi{\Alph{enumi})}
\tightlist
\item
  Il sistema formatta automaticamente la memoria per prevenire danni
\item
  L'attaccante può deviare il flusso del programma verso un indirizzo di
  sua scelta
\item
  Il router isola l'host compromesso dalla rete
\item
  Il programma viene ricompilato in modalità sicura
\end{enumerate}

\begin{enumerate}
\def\labelenumi{\arabic{enumi}.}
\setcounter{enumi}{8}
\tightlist
\item
  In ambito di Network Security, quale intervento rappresenta una
  ``mitigazione operativa'' (seppur non scalabile) contro l'ARP
  Spoofing?
\end{enumerate}

\begin{enumerate}
\def\labelenumi{\Alph{enumi})}
\tightlist
\item
  L'uso di cavi schermati
\item
  L'impostazione manuale di Tabelle ARP Statiche sugli host critici
\item
  Il riavvio periodico dello switch
\item
  La disabilitazione della porta 80
\end{enumerate}

Nome e Cognome:
\_\_\_\_\_\_\_\_\_\_\_\_\_\_\_\_\_\_\_\_\_\_\_\_\_\_\_\_\_\_\_\_\_\_\_\_\_\_\_\_\_\_\_\_\_\_\_\_\_\_\_\_\_
Data: \_\_\_\_\_\_\_\_\_\_\_\_\_\_\_\_\_\_\_

Punteggio: \_\_\_\_\_/20

\begin{enumerate}
\def\labelenumi{\arabic{enumi}.}
\setcounter{enumi}{9}
\tightlist
\item
  In ambito di Network Security, in cosa cosa consiste la Dynamic ARP
  Inspection? ?
\end{enumerate}

\begin{enumerate}
\def\labelenumi{\Alph{enumi})}
\tightlist
\item
  Bloccare l'ARP per permettere alla rete di funzionare solo a livello 3
\item
  Riconoscere il MAC address del gateway dinamicamente
\item
  Utilizzo di richieste DHCP snooping per mantenere associazione tra
  MAC-Address e IP address
\item
  Abilitare la comunicazione tra host wireless e wired
\end{enumerate}

\begin{enumerate}
\def\labelenumi{\arabic{enumi}.}
\setcounter{enumi}{10}
\tightlist
\item
  Cos'è un attacco di tipo ``Rogue DHCP''?
\end{enumerate}

\begin{enumerate}
\def\labelenumi{\Alph{enumi})}
\tightlist
\item
  Un server DHCP legittimo che esaurisce improvvisamente gli indirizzi
  IP disponibili
\item
  Un malware che disabilita il servizio DHCP sul router principale
  dell'azienda
\item
  L'introduzione nella rete di un server DHCP non autorizzato, che
  assegna configurazioni IP malevole (es. gateway e DNS errati) agli
  ignari client
\item
  Una tecnica crittografica per nascondere il traffico DHCP agli sniffer
  di rete
\end{enumerate}

\begin{enumerate}
\def\labelenumi{\arabic{enumi}.}
\setcounter{enumi}{11}
\tightlist
\item
  Qual è l'obiettivo primario di un attacco di ``DHCP Starvation''?
\end{enumerate}

\begin{enumerate}
\def\labelenumi{\Alph{enumi})}
\tightlist
\item
  Bloccare la porta 80 del server web aziendale
\item
  Esaurire tutti gli indirizzi IP disponibili nel pool del server DHCP
  legittimo inviando continue e false richieste di assegnazione
\item
  Rubare le password di amministratore memorizzate nel server DHCP
\item
  Forzare il server DHCP a fare il downgrade del protocollo alla
  versione precedente
\end{enumerate}

\begin{enumerate}
\def\labelenumi{\arabic{enumi}.}
\setcounter{enumi}{12}
\tightlist
\item
  Come funziona tipicamente un attacco di ``DNS Spoofing'' (o DNS
  Poisoning) all'interno di una rete locale?
\end{enumerate}

\begin{enumerate}
\def\labelenumi{\Alph{enumi})}
\tightlist
\item
  Inviando milioni di pacchetti ICMP (ping) al server DNS per farlo
  crashare
\item
  Sfruttando un Buffer Overflow nel browser della vittima
\item
  Rispondendo alle richieste DNS della vittima con indirizzi IP
  falsificati, reindirizzando così il suo traffico web verso un server
  controllato dall'attaccante
\item
  Cancellando fisicamente il file hosts sul computer del bersaglio
\end{enumerate}

\begin{enumerate}
\def\labelenumi{\arabic{enumi}.}
\setcounter{enumi}{13}
\tightlist
\item
  Durante l'analisi di un'applicazione web, qual è lo scopo principale
  dell'utilizzo di un tool come FFUF per la ``Directory e Domain
  Enumeration''?
\end{enumerate}

\begin{enumerate}
\def\labelenumi{\Alph{enumi})}
\tightlist
\item
  Decifrare le password degli utenti tramite un attacco a forza bruta
  sul form di login
\end{enumerate}

Nome e Cognome:
\_\_\_\_\_\_\_\_\_\_\_\_\_\_\_\_\_\_\_\_\_\_\_\_\_\_\_\_\_\_\_\_\_\_\_\_\_\_\_\_\_\_\_\_\_\_\_\_\_\_\_\_\_
Data: \_\_\_\_\_\_\_\_\_\_\_\_\_\_\_\_\_\_\_

Punteggio: \_\_\_\_\_/20

\begin{enumerate}
\def\labelenumi{\Alph{enumi})}
\setcounter{enumi}{1}
\tightlist
\item
  Scoprire cartelle, file nascosti o sottodomini non pubblici sul server
  web inviando rapidamente migliaia di richieste automatizzate basate su
  un dizionario
\item
  Intercettare in chiaro il traffico di rete tra il client e il server
  web
\item
  Eseguire comandi di amministrazione direttamente sul database del sito
\end{enumerate}

\begin{enumerate}
\def\labelenumi{\arabic{enumi}.}
\setcounter{enumi}{14}
\tightlist
\item
  Quale delle seguenti definizioni descrive meglio una vulnerabilità di
  ``SQL Injection'' (SQLi)?
\end{enumerate}

\begin{enumerate}
\def\labelenumi{\Alph{enumi})}
\tightlist
\item
  L'inserimento di input utente non filtrato all'interno di una query
  diretta al database, che permette all'attaccante di manipolare la
  richiesta o estrarre dati sensibili
\item
  L'inserimento di codice JavaScript maligno che verrà eseguito nel
  browser degli altri visitatori del sito
\item
  Un difetto di configurazione del server web che permette l'accesso
  alle cartelle senza password
\item
  Un attacco DoS (Denial of Service) che satura la banda di rete del
  server database
\end{enumerate}

\begin{enumerate}
\def\labelenumi{\arabic{enumi}.}
\setcounter{enumi}{15}
\tightlist
\item
  In ambito di Penetration Testing, qual è la caratteristica distintiva
  di una ``Reverse Shell'' rispetto a una shell tradizionale (Bind
  Shell)?
\end{enumerate}

\begin{enumerate}
\def\labelenumi{\Alph{enumi})}
\tightlist
\item
  La Reverse Shell è un'interfaccia grafica, mentre la Bind Shell è solo
  a riga di comando
\item
  La Reverse Shell si può utilizzare solo su sistemi operativi Linux
\item
  Nella Reverse Shell è il sistema vittima (compromesso) che avvia
  attivamente la connessione verso l'indirizzo IP della macchina
  dell'attaccante in ascolto
\item
  La Reverse Shell richiede sempre l'inserimento di una password di root
\end{enumerate}

\begin{enumerate}
\def\labelenumi{\arabic{enumi}.}
\setcounter{enumi}{16}
\tightlist
\item
  In un MAC Address, a cosa servono i primi 24 bit (chiamati OUI)?
\end{enumerate}

\begin{enumerate}
\def\labelenumi{\Alph{enumi})}
\tightlist
\item
  A identificare univocamente la singola scheda di rete prodotta
\item
  A stabilire la classe dell'indirizzo IP associato
\item
  A identificare il vendor (produttore) dell'interfaccia di rete
\item
  A determinare la priorità del traffico VLAN
\end{enumerate}

\begin{enumerate}
\def\labelenumi{\arabic{enumi}.}
\setcounter{enumi}{17}
\tightlist
\item
  Nella segmentazione della memoria di un processo, cosa contiene la
  sezione ``Text'' (o Code)?
\end{enumerate}

\begin{enumerate}
\def\labelenumi{\Alph{enumi})}
\tightlist
\item
  I file di testo aperti dall'utente
\item
  L'allocazione dinamica degli oggetti
\item
  Le istruzioni macchina eseguibili del programma, tipicamente in sola
  lettura
\item
  I messaggi di errore del compilatore
\end{enumerate}

Nome e Cognome:
\_\_\_\_\_\_\_\_\_\_\_\_\_\_\_\_\_\_\_\_\_\_\_\_\_\_\_\_\_\_\_\_\_\_\_\_\_\_\_\_\_\_\_\_\_\_\_\_\_\_\_\_\_
Data: \_\_\_\_\_\_\_\_\_\_\_\_\_\_\_\_\_\_\_

Punteggio: \_\_\_\_\_/20

\begin{enumerate}
\def\labelenumi{\arabic{enumi}.}
\setcounter{enumi}{18}
\tightlist
\item
  Qual è lo scopo primario del protocollo ARP (Address Resolution
  Protocol)?
\end{enumerate}

\begin{enumerate}
\def\labelenumi{\Alph{enumi})}
\tightlist
\item
  Risolvere i nomi a dominio nel corrispondente indirizzo IP
\item
  Assegnare un indirizzo IP dinamico a un nuovo host
\item
  Crittografare la comunicazione locale tra due switch
\item
  Conoscere il MAC Address di una macchina di cui si conosce già
  l'indirizzo
\end{enumerate}

\begin{enumerate}
\def\labelenumi{\arabic{enumi}.}
\setcounter{enumi}{19}
\tightlist
\item
  Nei processori con architettura x86 a 32-bit, qual è il registro
  critico che contiene l'indirizzo di memoria della prossima istruzione
  che la CPU dovrà eseguire?
\end{enumerate}

\begin{enumerate}
\def\labelenumi{\Alph{enumi})}
\tightlist
\item
  EAX
\item
  EBP
\item
  ESP
\item
  EIP
\end{enumerate}

N. Domanda

Risposta Corretta

Argomento Principale / Parola Chiave

1

C

Obiettivo Binary Exploitation (Arbitrary Code Execution)

2

D

Registro EIP (Instruction Pointer a 32-bit)

3

B

Debolezza del C (Mancanza di bounds checking)

4

C

Dimensione del MAC Address (48 bit)

5

C

Comportamento Switch (Flooding su Unknown Unicast)

6

A

Vulnerabilità ARP (Assenza totale di autenticazione)

7

C

Attacco ARP

8

B

Return address overwriting

9

B

ARP spoofing mitigazione operativa

10

C

Definizione DAI

11

C

DHCP Rogue definizione

12

B

DHCP startvation definizione

13

C

DNS spoofing definizione

14

B

Domain enumeration con fuzzing

15

A

SQL injection definizione

16

C

Definizione reverse shell

17

C

Struttura del MAC address

18

C

Cosa fa la sezione TEXT

19

D

ARP Protocol

20

D

Registri X86



\end{document}
